\documentclass[aspectratio=169,10pt]{beamer}

\usetheme{default}
\setbeamertemplate{navigation symbols}{}
\setbeamertemplate{footline}{}

\usepackage{booktabs}
\usepackage{array}
\usepackage{graphicx}
\usepackage{tikz}

\definecolor{Ink}{RGB}{28,37,44}
\definecolor{Slate}{RGB}{84,99,109}
\definecolor{Sand}{RGB}{245,240,232}
\definecolor{Teal}{RGB}{25,110,115}
\definecolor{Rust}{RGB}{179,91,45}
\definecolor{Gold}{RGB}{194,148,70}

\setbeamercolor{background canvas}{bg=Sand}
\setbeamercolor{normal text}{fg=Ink,bg=Sand}
\setbeamercolor{frametitle}{fg=Ink,bg=Sand}
\setbeamercolor{title}{fg=Ink}
\setbeamercolor{block title}{fg=Sand,bg=Teal}
\setbeamercolor{block body}{fg=Ink,bg=white}
\setbeamercolor{block title alerted}{fg=Sand,bg=Rust}
\setbeamercolor{block body alerted}{fg=Ink,bg=white}
\setbeamercolor{itemize item}{fg=Rust}
\setbeamercolor{itemize subitem}{fg=Rust}

\setbeamerfont{title}{series=\bfseries,size=\Large}
\setbeamerfont{frametitle}{series=\bfseries,size=\large}
\setbeamerfont{block title}{series=\bfseries}

\newcommand{\topband}{
    \begin{tikzpicture}[remember picture,overlay]
        \fill[Teal] (current page.north west) rectangle ([yshift=-0.34cm]current page.north east);
        \fill[Gold] ([yshift=-0.34cm]current page.north west) rectangle ([xshift=2.3cm,yshift=-0.48cm]current page.north west);
    \end{tikzpicture}
}

\title{QIZ Model: Motivation, Design, and Early Results}
\author{}
\date{}

\begin{document}

\begin{frame}[t]
\topband
\vspace{0.35cm}
\frametitle{Why a Model, and Why This One?}

\begin{columns}[T,onlytextwidth]
    \begin{column}{0.41\textwidth}
        \begin{block}{Why a model is needed}
        \small
        \begin{itemize}
            \item QIZ changes several margins at once: tariff-free US access, compliance costs, sourcing requirements, export participation, and firm upgrading.
            \item The policy can help some firms and regions while hurting others, so reduced-form intuition alone is not enough.
            \item I need a framework that links micro choices by firms to aggregate outcomes such as wages, labor reallocation, exports, and welfare.
        \end{itemize}
        \end{block}

        \begin{alertblock}{Why this model is suitable}
        \small
        \begin{itemize}
            \item Heterogeneous firms are essential because only some firms export, comply, or upgrade.
            \item General equilibrium is essential because QIZ affects wages, price indices, entry, and movement of labor across Egyptian regions.
            \item The model is disciplined enough for counterfactuals: QIZ on/off, same-type complier comparisons, and changes in ROO stringency.
        \end{itemize}
        \end{alertblock}
    \end{column}

    \begin{column}{0.56\textwidth}
        \begin{block}{Model structure in the current version}
        \centering
        \resizebox{\linewidth}{!}{\begin{tikzpicture}[
    font=\scriptsize,
    >=Latex,
    node distance=0.35cm and 0.45cm,
    box/.style={draw=Teal, very thick, rounded corners=2mm, fill=white, align=left},
    soft/.style={draw=Slate, semithick, rounded corners=2mm, fill=white, align=left},
    accent/.style={draw=Rust, very thick, rounded corners=2mm, fill=white, align=left},
    flow/.style={-Latex, thick, draw=Ink},
    policy/.style={-Latex, very thick, draw=Rust},
    feedback/.style={-Latex, semithick, draw=Gold!70!Ink}
]

\node[soft, minimum width=3.15cm, minimum height=0.95cm] (types) at (0,3.2)
{\textbf{Firm heterogeneity}\\productivity $\phi$ \\ compliance cost $\varepsilon$};

\node[box, minimum width=5.0cm, minimum height=2.9cm] (egypt) at (0,0.85) {};
\node[anchor=north west, font=\bfseries\small, text=Teal] at ([xshift=0.12cm,yshift=-0.1cm]egypt.north west)
{Egyptian production side};

\node[soft, minimum width=2.05cm, minimum height=1.15cm] (q) at (-1.2,0.95)
{\textbf{Q region}\\QIZ-eligible\\can comply};
\node[soft, minimum width=2.05cm, minimum height=1.15cm] (n) at (1.2,0.95)
{\textbf{N region}\\non-QIZ\\no preference};

\node[soft, minimum width=4.25cm, minimum height=0.82cm] (sectors) at (0,-0.32)
{\textbf{Sectors now:} Textiles ($T$) and Other ($O$) \hfill expand later};

\node[accent, minimum width=2.95cm, minimum height=1.35cm] (choices) at (4.4,1.0)
{\textbf{Firm choices}\\serve EG / US / RW\\comply with ROO?\\upgrade productivity?};

\node[soft, minimum width=1.5cm, minimum height=0.72cm] (eg) at (8.05,2.0) {\textbf{Egypt}\\final sales};
\node[soft, minimum width=1.5cm, minimum height=0.72cm] (us) at (8.05,1.0) {\textbf{US}\\final sales};
\node[soft, minimum width=1.5cm, minimum height=0.72cm] (rw) at (8.05,0.0) {\textbf{RW}\\final sales};

\node[soft, minimum width=1.95cm, minimum height=0.92cm] (israel) at (5.9,2.55)
{\textbf{Israel}\\input source for\\ROO content};

\node[soft, minimum width=5.2cm, minimum height=0.95cm] (ge) at (2.55,-1.48)
{\textbf{GE feedback:} wages, price indices, entry, labor mobility, welfare};

\draw[flow] (types.south) -- ([yshift=0.05cm]egypt.north);
\draw[flow] (q.east) -- (choices.west);
\draw[flow] (n.east) -- (choices.west);
\draw[flow] (choices.east) -- (eg.west);
\draw[policy] (choices.east) -- node[above, sloped, text=Rust] {\tiny tariff-free if compliant} (us.west);
\draw[flow] (choices.east) -- (rw.west);
\draw[policy] (israel.south west) -- node[right, text=Rust] {\tiny ROO input share $\gamma$} ([xshift=0.1cm,yshift=0.12cm]choices.north west);
\draw[feedback] (ge.north west) to[out=105,in=-90] (q.south west);
\draw[feedback] (ge.north) to[out=85,in=-90] (n.south);
\draw[feedback] (ge.north east) to[out=75,in=-90] (choices.south east);

\node[anchor=west, text=Slate, font=\tiny] at (-2.45,-2.15)
{Four relevant locations economically: Egypt, US, Israel, RW.};
\node[anchor=west, text=Slate, font=\tiny] at (-2.45,-2.45)
{Three final-demand markets in the current code: EG, US, RW.};

\end{tikzpicture}
}
        \end{block}
    \end{column}
\end{columns}
\end{frame}

\begin{frame}[t]
\topband
\vspace{0.35cm}
\frametitle{What the Model Delivers So Far}

\begin{columns}[T,onlytextwidth]
    \begin{column}{0.54\textwidth}
        \begin{block}{Benchmark comparison: QIZ on vs no QIZ}
        \scriptsize
        \begin{tabular}{>{\raggedright\arraybackslash}p{3.15cm}ccc}
        \toprule
        Outcome & No QIZ & QIZ & Change \\
        \midrule
        Welfare & 127.45 & 131.94 & +3.52\% \\
        Real wage in Q & 1.315 & 1.371 & +4.25\% \\
        Real wage in N & 1.227 & 1.258 & +2.46\% \\
        Total US exports & 45.93 & 51.67 & +12.49\% \\
        Q employment share & 53.46\% & 54.32\% & +0.86 pp \\
        \bottomrule
        \end{tabular}
        \end{block}

        \begin{block}{Firm-level mechanism in Q-textiles}
        \small
        Same-type firms that comply under QIZ increase:
        \begin{itemize}
            \item US exports by \textbf{376.7\%}
            \item non-US exports by \textbf{129.2\%}
        \end{itemize}
        In the current benchmark, \textbf{50\%} of active Q-textile firms comply.
        \end{block}
    \end{column}

    \begin{column}{0.42\textwidth}
        \begin{alertblock}{Interpretation}
        \small
        The current calibration suggests that QIZ improves aggregate performance mainly through better US market access, higher real wages, and a modest reallocation of labor toward QIZ regions.
        \end{alertblock}

        \begin{block}{ROO sensitivity}
        \small
        In a separate gamma counterfactual, increasing the Israeli-input requirement from \textbf{0\%} to \textbf{30\%}:
        \begin{itemize}
            \item lowers compliance from \textbf{64.9\%} to \textbf{12.5\%}
            \item lowers Q-to-US exports from \textbf{18.0} to \textbf{14.1}
            \item lowers welfare from \textbf{126.90} to \textbf{125.04}
        \end{itemize}
        This is consistent with stricter ROO making preferential access less valuable.
        \end{block}

        \begin{block}{Where the project stands}
        \small
        These are \textbf{preliminary calibrated results}, not structural estimates yet. The next extension is finer sectoral disaggregation beyond textiles vs other manufacturing.
        \end{block}
    \end{column}
\end{columns}
\end{frame}

\end{document}
